\documentclass[10pt,twocolumn,letterpaper]{article}

\usepackage[review]{cvpr}
%% This file contains a number of tweaks that are typically applied to the main document.
%% They are not enabled by default, but can be enabled by uncommenting the relevant lines.


%%
%% Official support for 2-column teaser figure
%%
\usepackage{cuted} %< for strip environment

%%
%% For auto-referencing of figures (see fig/teaser.tex and \label{} therein) 
%%
\usepackage{currfile} %< for \currfilebase macro

%%
%% caption and (global) spacing overrides
%%
\usepackage{caption} %< for "\captionof" commands (teaser)
% \setlength{\abovecaptionskip}{.5em} % space above caption
% \setlength{\belowcaptionskip}{1em} % space below caption

%%
%% Inline annotations; for predefined colors, refer to "dvipsnames" in the xcolor package:
%% https://tinyurl.com/overleaf-colors
%%
\newcommand{\red}[1]{{\color{red}#1}}
\newcommand{\todo}[1]{{\color{red}#1}}
\newcommand{\TODO}[1]{\textbf{\color{red}[TODO: #1]}}
%%
%% Disable for camera ready / submission by uncommenting these lines:
%%
% \renewcommand{\TODO}[1]{}
% \renewcommand{\todo}[1]{#1}

%%
%% Work harder in optimizing text layout. Typically shrinks text by 1/6 of page, enable
%% it at the very end of the writing process, when you are just above the page limit.
%%
% \usepackage{microtype}

%%
%% Fine-tune paragraph spacing.
%%
% \renewcommand{\paragraph}[1]{\vspace{.5em}\noindent\textbf{#1.}}

%%
%% Globally adjusts space between figure and caption.
%%
% \setlength{\abovecaptionskip}{.5em}

%%
%% Allows "the use of \paper to refer to the project name"
%% with automatic management of space at the end of the word.
%%
% \usepackage{xspace}
% \newcommand{\paper}{ProjectName\xspace}

%%
%% Commonly used math definitions.
%%
% \DeclareMathOperator*{\argmin}{arg\,min}
% \DeclareMathOperator*{\argmax}{arg\,max}

%%
%% Tighten underline.
%%
% \usepackage{soul}
% \setuldepth{foobar}

%%
%% Reconfigure paragraph (tweaks spacing)
%%

\usepackage{xr}
\externaldocument{main}
\usepackage{booktabs}
\usepackage{multirow}
\usepackage{graphicx}
\usepackage{amsmath,amssymb,mathtools}

\definecolor{cvprblue}{rgb}{0.21,0.49,0.74}
\usepackage[pagebackref,breaklinks,colorlinks,allcolors=cvprblue]{hyperref}

\def\paperID{*****}
\def\confName{CVPR}
\def\confYear{2027}

% same generated numbers as the main paper
% placeholder: run scripts/make_paper_numbers.py after the full Kaggle results are in outputs/results/

\newcommand{\maybefig}[1]{\IfFileExists{fig/#1.pdf}{\includegraphics[width=\linewidth]{fig/#1.pdf}}{\fbox{\small figure #1 pending}}}
\renewcommand{\thesection}{\Alph{section}}
% S-prefixed numbers, so they cannot be mistaken for the main paper's
\renewcommand{\thetable}{S\arabic{table}}
\renewcommand{\thefigure}{S\arabic{figure}}

\title{Zoom In, Scale Out: Metric Depth Models Do Not Respect a Digital Zoom\\[2pt]
\large Supplementary Material}

\author{Anonymous CVPR submission}

\begin{document}
\maketitle
\sloppy

Section and table numbers without a letter refer to the main paper. Every number here is generated
from the same result files as the main paper (\texttt{scripts/}\allowbreak\texttt{make\_paper\_numbers.py}).

\section{Qualitative example on NYUv2}

\begin{figure}[h]
  \centering
  \maybefig{figS_teaser_nyu}
  \caption{The layout of \cref{fig:teaser} on an NYUv2 test image: depth at $1\times$, $2\times$ and
  $3\times$ digital zoom on one colour scale in metres. The number in each panel is the median depth over
  the pixels inside the $3\times$ crop, the same physical points in every column.}
  \label{fig:supp_nyu}
\end{figure}

\section{The depth exponent per dataset}

\Cref{tab:alpha_ds} and \cref{fig:supp_alpha_ds} split the pooled $\alpha$ of \cref{tab:audit} by
dataset. On every set each pixels-only model has a larger $\alpha$ than each camera-predicting model
(ZoeDepth and MoGe-2-S tie on DIODE-indoor), and Metric3D v2 with the true camera stays near zero. The
pixels-only models shrink most on KITTI and NYUv2, where DA-V2-S even exceeds pure apparent-size
behaviour ($\alpha>1$).

\begin{table}[t]\centering\small
\caption{Depth exponent $\alpha$ per dataset (median over images).}\label{tab:alpha_ds}
\setlength{\tabcolsep}{3pt}
\resizebox{\linewidth}{!}{\begin{tabular}{lcccc}\toprule
Model & KITTI & NYUv2 & DIODE-in & DIODE-out \\ \midrule
DA-V2-S & 1.28 & 1.07 & 0.76 & 0.71 \\
DA-V2-B & 1.19 & 0.83 & 0.75 & 0.66 \\
ZoeDepth & 0.82 & 0.60 & 0.43 & 0.52 \\
Depth Pro & 0.35 & 0.46 & 0.19 & -0.01 \\
MoGe-2-S & 0.44 & 0.21 & 0.43 & 0.33 \\
UniDepthV2-S & 0.60 & 0.13 & 0.32 & 0.25 \\
Metric3D v2-S & -0.03 & -0.11 & -0.09 & -0.18 \\
Metric3D v2-S (stale K) & 0.97 & 0.89 & 0.91 & 0.82 \\
\bottomrule\end{tabular}}\end{table}


\begin{figure}[t]
  \centering
  \maybefig{figS_alpha_per_dataset}
  \caption{Median $\alpha$ per model and dataset.}
  \label{fig:supp_alpha_ds}
\end{figure}

\section{Off-centre crops and the principal point}

Zooms about the four quadrant centres (50 images per set from KITTI and NYUv2) give $\alpha$ of the same
order as centred zooms (\cref{fig:supp_offc}). For such crops the zoom law moves the principal point
by $(c_x-x_0+\tfrac12)\,W/w-\tfrac12$, half the image width at $2\times$. \Cref{fig:supp_pp} compares the
principal point predicted by MoGe-2 and UniDepthV2 with the law: both errors lie on the curve of a
predictor that always answers the image centre (\MogeTwoSPpQuadTwo{} and \UnidepthTwoSPpQuadTwo{} of
the width at $2\times$, against \MogeTwoSPpGuessTwo).

\begin{figure}[t]
  \centering
  \maybefig{figS_offcentre_alpha}
  \caption{Median $\alpha$ by crop position.}
  \label{fig:supp_offc}
\end{figure}

\begin{figure}[t]
  \centering
  \maybefig{figS_principal_point}
  \caption{Principal-point error (fraction of the image width) under quadrant zooms. Dotted: the error
  of always answering the image centre.}
  \label{fig:supp_pp}
\end{figure}

\section{ZoomFix per dataset}

\Cref{tab:zoomfix} gives the cross-fitted ZoomFix of \cref{sec:results} per dataset. The label-free
exponent tracks the dataset: it is above 1 for DA-V2-S on KITTI and NYUv2, where the pooled $\alpha$ hides
a stronger shrink. The MoGe-2-S entry for DIODE-outdoor is dominated by a few zoomed images with very
large errors (the zoomed AbsRel is \MogeTwoSFixDiodeOutZoomed{} against \MogeTwoSFixDiodeOutTarget{}
unzoomed), which a global factor cannot repair.

\begin{table*}[t]\centering\scriptsize
\caption{ZoomFix per dataset at $3\times$, cross-fitted: the exponent $\hat e$ is measured without labels on half of the images and the correction scored on the other half. AbsRel of the zoomed prediction as is, after multiplying by $s^{\hat e}$, and of the unzoomed prediction on the same pixels (the target).}\label{tab:zoomfix}
\setlength{\tabcolsep}{3pt}
\begin{tabular}{l cccc cccc cccc cccc}\toprule
& \multicolumn{4}{c}{KITTI} & \multicolumn{4}{c}{NYUv2} & \multicolumn{4}{c}{DIODE-in} & \multicolumn{4}{c}{DIODE-out} \\
\cmidrule(lr){2-5}\cmidrule(lr){6-9}\cmidrule(lr){10-13}\cmidrule(lr){14-17}
Model & $\hat e$ & as is & fixed & target & $\hat e$ & as is & fixed & target & $\hat e$ & as is & fixed & target & $\hat e$ & as is & fixed & target \\ \midrule
DA-V2-S & 1.25 & 0.758 & 0.384 & 0.149 & 1.12 & 0.675 & 0.495 & 0.205 & 0.75 & 0.598 & 0.365 & 0.263 & 0.69 & 0.558 & 0.661 & 0.605 \\
Depth Pro & 0.31 & 0.258 & 0.277 & 0.160 & 0.44 & 0.449 & 0.223 & 0.121 & 0.19 & 0.399 & 0.421 & 0.201 & 0.12 & 0.636 & 0.638 & 0.665 \\
MoGe-2-S & 0.44 & 0.536 & 0.262 & 0.255 & 0.12 & 0.272 & 0.227 & 0.129 & 0.44 & 0.493 & 0.364 & 0.257 & 0.38 & 3.037 & 4.396 & 0.379 \\
\bottomrule\end{tabular}\end{table*}


\section{Fine-tuning details}

\paragraph{Variants.} All variants start from the same DA-V2-S checkpoint and differ only in the loss:
EFT (anchor + zoom term), EFT+B (+ blur term), EFT+B with $\lambda=3$, blur-only (anchor + blur term),
and the Zhong \etal-style target (anchor + zoom term, with the student's own stop-gradient prediction on
the unzoomed image as target). Each is trained once per domain (KITTI raw outdoor, NYUv2 train indoor)
with the settings of \cref{sec:eft}.

\paragraph{Training curves.} \Cref{fig:supp_train} shows the zoom and anchor terms over training
(running mean over 100 steps). The anchor term is exactly zero at the first step, since the student is a
copy of the teacher; it jumps as soon as the zoom term moves the weights and then falls through training
(EFT outdoor: mean 0.10 over the first 100 steps, 0.04 over the last 100), as the model learns to
satisfy both terms.

\begin{figure}[t]
  \centering
  \maybefig{figS_eft_training}
  \caption{Zoom loss (solid) and anchor loss (dashed) of the fine-tuning variants, outdoor and indoor
  heads. Blur-only has no zoom term.}
  \label{fig:supp_train}
\end{figure}

\section{Models and compute}

\paragraph{Checkpoints.} Hugging Face ids unless noted.
{\footnotesize\raggedright
\begin{itemize}\setlength{\itemsep}{0pt}
\item DA-V2: \texttt{depth-anything/Depth-Anything-V2-} \texttt{Metric-\{Indoor,Outdoor\}-\{Small,Base,Large\}-hf}
\item ZoeDepth: \texttt{Intel/zoedepth-nyu-kitti}
\item Depth Pro: \texttt{apple/DepthPro-hf}
\item MoGe-2: \texttt{Ruicheng/moge-2-\{vits,vitb,vitl\}-normal}
\item UniDepthV2: \texttt{lpiccinelli/unidepth-v2-} \texttt{\{vits14,vitb14,vitl14\}}
\item UniK3D: \texttt{lpiccinelli/unik3d-vitl}
\item DA3: \texttt{depth-anything/DA3NESTED-GIANT-} \texttt{LARGE-1.1}, \texttt{DA3METRIC-LARGE}, \texttt{DA3-LARGE-1.1}
\item VGGT: \texttt{facebook/VGGT-1B}
\item Metric3D v2: \texttt{metric3d\_vit\_\{small,large\}} (official torch hub entry)
\end{itemize}}

\paragraph{Compute.} All audits and fine-tuning runs ran as Kaggle notebooks with one GPU each; the
analysis runs on a laptop CPU. The code release contains the kernel builders, the run logs and every
result CSV, and \texttt{scripts/verify\_results.py} checks that all models saw the same images and all
zoom rows are complete.

\end{document}
